% ---------------------------------------------------------------
% Evaluation. Holds everything about MEASUREMENT; the experimental
% setup section holds data, training and the systems compared, and
% deliberately says nothing here (2026-09-13, by request).
%
% Moved here from experiments.tex, unedited, as a starting point:
% \subsection{Protocol and metrics}
% \label{sec:protocol}
%% Each system receives a five-bar prompt (80 frames) and generates to 416
% frames, 21 bars of continuation, with three samples per song at
% temperature $1$. Metrics are computed on the generated span only and
% are calibrated against the ground-truth continuation of the same song,
% so a system is measured against what that song actually did rather than
% against a corpus average.
%% We report three families. \emph{Coupling}: chord-tone coverage, the
% share of generated melody notes whose pitch class lies in the
% concurrent generated chord, reported as a delta against the same
% statistic on the reference pair. \emph{Distributional fit}: Jensen--
% Shannon divergences between generated and reference distributions of
% harmonic rhythm, pitch class, note density and duration.
% \emph{Survival}: the fraction of frames left empty, which catches the
% failure where one stream falls silent and the other continues alone.
% A listening test [X] covers what the objective metrics cannot.
% ---------------------------------------------------------------
\section{Evaluation}
\label{sec:evaluation}

% TODO
